% Klausur 1 (29.07.2026): Lösungen. Alle Zahlen mit den Werkzeugen in skripte/ berechnet und geprüft.
% Gemeinsame Präambel für alle LaTeX-Dateien.
\documentclass[a4paper,11pt]{article}
\usepackage[utf8]{inputenc}
\usepackage[T1]{fontenc}
\usepackage[ngerman]{babel}
\usepackage{amsmath,amssymb,amsthm}
\usepackage{enumitem}
\usepackage{booktabs}
\setcounter{MaxMatrixCols}{30}

\theoremstyle{definition}
\newtheorem{definition}{Definition}[section]
\newtheorem{satz}[definition]{Satz}
\newtheorem{lemma}[definition]{Lemma}
\newtheorem{korollar}[definition]{Korollar}
\newtheorem{beispiel}[definition]{Beispiel}
\newtheorem{bemerkung}[definition]{Bemerkung}
\newtheorem{algorithmus}[definition]{Algorithmus}
\newtheorem{aufgabe}{Aufgabe}[section]
\newenvironment{loesung}{\par\noindent\textbf{Lösung.}\ }{\par}

% Kurzbefehle
\newcommand{\N}{\mathbb{N}}
\newcommand{\Z}{\mathbb{Z}}
\newcommand{\Q}{\mathbb{Q}}
\newcommand{\R}{\mathbb{R}}
\newcommand{\Zm}[1]{\mathbb{Z}_{#1}}
\newcommand{\ggT}{\operatorname{ggT}}
\newcommand{\kgV}{\operatorname{kgV}}

\usepackage[a4paper,margin=2.2cm]{geometry}
\usepackage{fancyhdr}
\usepackage{xcolor}
\usepackage[most]{tcolorbox}
\pagestyle{fancy}\fancyhf{}
\lhead{Diskrete Mathematik -- Lösungen Klausur 29.07.2026}\rhead{Seite \thepage}
\setlength{\parindent}{0pt}\setlength{\parskip}{4pt}
\setcounter{secnumdepth}{1}

\newtcolorbox{ergebnis}{colback=green!6,colframe=green!45!black,boxrule=0.6pt,left=6pt,right=6pt,top=3pt,bottom=3pt,sharp corners}
\newtcolorbox{achtung}{colback=orange!8,colframe=orange!70!black,boxrule=0.6pt,left=6pt,right=6pt,top=3pt,bottom=3pt,sharp corners}
\newcommand{\teil}[1]{\par\medskip\textbf{#1)}\ }

\begin{document}

\begin{center}
{\LARGE\bfseries Lösungen zur Klausur Diskrete Mathematik}\\[4pt]
{\large 29.07.2026}
\end{center}

\begin{center}
\begin{tabular}{@{}lll@{}}
\toprule
Aufgabe & Thema & Ergebnis\\
\midrule
1 & RSA & $e = 11011101_2$, $c = 896$, $\varphi(m) = 1512$, $d = 821$\\
2 & Inklusion -- Exklusion & $288$\\
3 & Euklidischer Algorithmus & $\ggT(17,15) = 1$, $x = -7$, $y = -8$ (oder $x = 8$, $y = 9$)\\
4 & Chinesischer Restsatz & $z \equiv 110 \pmod{255}$, kleinste positive Lösung $110$\\
5 & Modulo-11-Code & $c = 1150000546$; $r_1$ nicht dekodierbar; $r_2 \to 1150020733$\\
\bottomrule
\end{tabular}
\end{center}

% ---------------------------------------------------------------------------
\section*{Aufgabe 1: RSA}
Gegeben: $w = 1511$, $m = 43 \cdot 57 = 2451$, $e = 221$.

\teil{a} \textbf{$e$ im Binärsystem.}
\[
221 = 128 + 64 + 16 + 8 + 4 + 1 = 2^7 + 2^6 + 2^4 + 2^3 + 2^2 + 2^0
\quad\Rightarrow\quad e = 11011101_2 .
\]

\teil{b} \textbf{Verschlüsseln.} Formel (Alice): $c \equiv w^e \pmod m$, also $c \equiv 1511^{221} \pmod{2451}$.

Schnelles Potenzieren: $1511^{221} = 1511^{1} \cdot 1511^{4} \cdot 1511^{8} \cdot 1511^{16} \cdot 1511^{64} \cdot 1511^{128}$.

Quadrattabelle (mod 2451):
\[
\begin{array}{r@{\;\equiv\;}l@{\qquad}l}
1511^{1}   & 1511 & \leftarrow\\
1511^{2}   & 1511^2 = 2283121 = 931 \cdot 2451 + 1240 \equiv 1240 & \\
1511^{4}   & 1240^2 = 1537600 = 627 \cdot 2451 + 823 \equiv 823 & \leftarrow\\
1511^{8}   & 823^2 = 677329 = 276 \cdot 2451 + 853 \equiv 853 & \leftarrow\\
1511^{16}  & 853^2 = 727609 = 296 \cdot 2451 + 2113 \equiv 2113 & \leftarrow\\
1511^{32}  & 2113^2 = 4464769 = 1821 \cdot 2451 + 1498 \equiv 1498 & \\
1511^{64}  & 1498^2 = 2244004 = 915 \cdot 2451 + 1339 \equiv 1339 & \leftarrow\\
1511^{128} & 1339^2 = 1792921 = 731 \cdot 2451 + 1240 \equiv 1240 & \leftarrow
\end{array}
\]
Produkt Schritt für Schritt (mod 2451):
\begin{align*}
1511 \cdot 823 &= 1243553 \equiv 896\\
896 \cdot 853 &= 764288 \equiv 2027\\
2027 \cdot 2113 &= 4283051 \equiv 1154\\
1154 \cdot 1339 &= 1545206 \equiv 1076\\
1076 \cdot 1240 &= 1334240 \equiv 896
\end{align*}
\begin{ergebnis} $c = 1511^{221} \bmod 2451 = 896$ \end{ergebnis}

\teil{c} \textbf{Entschlüsseln.} Formel (Bob): $w \equiv c^d \pmod m$.
Das funktioniert, weil $e \cdot d \equiv 1 \pmod{\varphi(m)}$ und nach Euler
$c^d \equiv (w^e)^d = w^{ed} \equiv w \pmod m$.

\teil{d} \textbf{$\varphi(m)$.} Achtung: $57 = 3 \cdot 19$ ist keine Primzahl. Vollständige Zerlegung:
\[
2451 = 3 \cdot 19 \cdot 43 .
\]
Da $\varphi$ multiplikativ ist und $\varphi(p) = p - 1$ für Primzahlen:
\[
\varphi(2451) = \varphi(3)\,\varphi(19)\,\varphi(43) = 2 \cdot 18 \cdot 42 = 1512 .
\]
\begin{ergebnis} $\varphi(m) = 1512$ \end{ergebnis}
\begin{achtung}
\textbf{Falle:} Wer $\varphi(m) = (p-1)(q-1) = 42 \cdot 56 = 2352$ rechnet, erhält $d = 221^{-1} \bmod 2352 = 149$.
Das ist falsch: $896^{149} \bmod 2451 = 737 \neq 1511$. Die Formel $(p-1)(q-1)$ gilt nur für Primzahlen $p \neq q$.
\end{achtung}

\teil{e} \textbf{Privater Schlüssel $d$.}
Oscar zerlegt $m$ in Primfaktoren, berechnet $\varphi(m)$ und löst
\[
e \cdot d \equiv 1 \pmod{\varphi(m)}, \quad\text{also}\quad 221 \cdot d \equiv 1 \pmod{1512}.
\]

\emph{Weg 1: Euklidischer Algorithmus (Lemma von Bezout).}
\[
\begin{array}{r@{\;=\;}l}
1512 & 6 \cdot 221 + 186\\
221 & 1 \cdot 186 + 35\\
186 & 5 \cdot 35 + 11\\
35 & 3 \cdot 11 + 2\\
11 & 5 \cdot 2 + 1\\
2 & 2 \cdot 1 + 0
\end{array}
\qquad\Rightarrow\quad \ggT(1512, 221) = 1 .
\]
Matrix $Q = \prod \begin{pmatrix} q_i & 1\\ 1 & 0\end{pmatrix}$ mit $q = 6,1,5,3,5,2$:
\[
Q = \begin{pmatrix} 1512 & 691\\ 221 & 101 \end{pmatrix}, \qquad
\det Q = 1512 \cdot 101 - 221 \cdot 691 = 1 .
\]
Also $221 \cdot (-691) \equiv 1 \pmod{1512}$ und $d = -691 \equiv 821 \pmod{1512}$.

\emph{Weg 2: Satz von Euler.} Aus $e^{\varphi(\varphi(m))} \equiv 1 \pmod{\varphi(m)}$ folgt
$d \equiv e^{\varphi(\varphi(m)) - 1} \pmod{\varphi(m)}$.
Mit $1512 = 2^3 \cdot 3^3 \cdot 7$ ist $\varphi(1512) = 4 \cdot 18 \cdot 6 = 432$, also
\[
d \equiv 221^{431} \pmod{1512}.
\]
$431 = 110101111_2 = 256 + 128 + 32 + 8 + 4 + 2 + 1$. Quadrattabelle (mod 1512):
$221^{1} = 221$, $221^{2} \equiv 457$, $221^{4} \equiv 193$, $221^{8} \equiv 961$, $221^{16} \equiv 1201$,
$221^{32} \equiv 1465$, $221^{64} \equiv 697$, $221^{128} \equiv 457$, $221^{256} \equiv 193$.
\begin{align*}
221 \cdot 457 &\equiv 1205, & 1205 \cdot 193 &\equiv 1229, & 1229 \cdot 961 &\equiv 197,\\
197 \cdot 1465 &\equiv 1325, & 1325 \cdot 457 &\equiv 725, & 725 \cdot 193 &\equiv 821 .
\end{align*}
Kontrolle: $221 \cdot 821 = 181441 = 120 \cdot 1512 + 1 \equiv 1 \pmod{1512}$.
\begin{ergebnis} $d = 821$ \end{ergebnis}

\emph{Probe:} $896^{821} \bmod 2451$ mit $821 = 1100110101_2$:
Quadrate $896^{1} = 896$, $896^{4} \equiv 1240$, $896^{16} \equiv 853$, $896^{32} \equiv 2113$, $896^{256} \equiv 1240$, $896^{512} \equiv 823$;
Produkt $896 \cdot 1240 \equiv 737$, $\cdot 853 \equiv 1205$, $\cdot 2113 \equiv 2027$, $\cdot 1240 \equiv 1205$, $\cdot 823 \equiv 1511 = w$. \checkmark

% ---------------------------------------------------------------------------
\section*{Aufgabe 2: Inklusion -- Exklusion}
$\Omega = \{1 \le k \le 720\}$, $|\Omega| = 720$, $A_i = \{k \in \Omega \mid i \text{ teilt } k\}$ für $i = 3, 4, 5$.
Gesucht: $|\Omega \setminus (A_3 \cup A_4 \cup A_5)| = |\Omega| - \alpha_1 + \alpha_2 - \alpha_3$ mit
\begin{align*}
\alpha_1 &= |A_3| + |A_4| + |A_5| = \tfrac{720}{3} + \tfrac{720}{4} + \tfrac{720}{5} = 240 + 180 + 144 = 564,\\
\alpha_2 &= |A_3 \cap A_4| + |A_3 \cap A_5| + |A_4 \cap A_5| = \tfrac{720}{12} + \tfrac{720}{15} + \tfrac{720}{20} = 60 + 48 + 36 = 144,\\
\alpha_3 &= |A_3 \cap A_4 \cap A_5| = \tfrac{720}{60} = 12 .
\end{align*}
(Schnittmengen über das kgV: $\kgV(3,4) = 12$, $\kgV(3,5) = 15$, $\kgV(4,5) = 20$, $\kgV(3,4,5) = 60$.)
\[
720 - 564 + 144 - 12 = 288 .
\]
Kontrolle: $720 \cdot (1 - \tfrac13)(1 - \tfrac14)(1 - \tfrac15) = 720 \cdot \tfrac23 \cdot \tfrac34 \cdot \tfrac45 = 288$.
\begin{ergebnis} $288$ Zahlen $\le 720$ sind durch keine der Zahlen $3, 4, 5$ teilbar. \end{ergebnis}

% ---------------------------------------------------------------------------
\section*{Aufgabe 3: Euklidischer Algorithmus}
\teil{a}
\[
\begin{array}{r@{\;=\;}l}
17 & 1 \cdot 15 + 2\\
15 & 7 \cdot 2 + 1\\
2 & 2 \cdot 1 + 0
\end{array}
\qquad\Rightarrow\quad \ggT(17, 15) = 1 .
\]

\teil{b} EA rückwärts:
\begin{align*}
1 &= 15 - 7 \cdot 2 = 15 - 7 \cdot (17 - 15) = 8 \cdot 15 - 7 \cdot 17 .
\end{align*}
Also $17 \cdot (-7) - 15 \cdot (-8) = -119 + 120 = 1$.

(Mit der Matrix $Q = \begin{pmatrix}1&1\\1&0\end{pmatrix}\begin{pmatrix}7&1\\1&0\end{pmatrix}\begin{pmatrix}2&1\\1&0\end{pmatrix} = \begin{pmatrix}17&8\\15&7\end{pmatrix}$,
$\det Q = 17 \cdot 7 - 15 \cdot 8 = -1$, ergibt sich dieselbe Lösung.)

Alle Lösungen: $x = -7 + 15t$, $y = -8 + 17t$ ($t \in \Z$). Für $t = 1$: $x = 8$, $y = 9$ ($17 \cdot 8 - 15 \cdot 9 = 136 - 135 = 1$).
\begin{ergebnis} $\ggT(17,15) = 1$; eine Lösung: $x = -7$, $y = -8$ (ebenso $x = 8$, $y = 9$). \end{ergebnis}

% ---------------------------------------------------------------------------
\section*{Aufgabe 4: Chinesischer Restsatz}
\teil{a} Ansatz: $z = 17x + 8 = 15y + 5$, also
\[
17x - 15y = 5 - 8 = -3 .
\]
Aus Aufgabe 3: $17 \cdot 8 - 15 \cdot 9 = 1$. Multiplikation mit $-3$:
\[
17 \cdot (-24) - 15 \cdot (-27) = -3, \qquad x = -24,\ y = -27 .
\]
\[
z = 17 \cdot (-24) + 8 = -400 = 15 \cdot (-27) + 5 .
\]
(Mit $x = -7$, $y = -8$ aus Aufgabe 3 erhält man $x = 21$, $y = 24$ und $z = 365$; das ist dieselbe Restklasse.)

\teil{b} Da $\ggT(17, 15) = 1$, sind alle Lösungen $z \equiv -400 \pmod{17 \cdot 15 = 255}$:
\[
z = -400 + k \cdot 255 \quad (k \in \Z).
\]
Kleinste positive Lösung: $z = -400 + 2 \cdot 255 = 110$.

Probe: $110 = 6 \cdot 17 + 8 \equiv 8 \pmod{17}$ \checkmark, \quad $110 = 7 \cdot 15 + 5 \equiv 5 \pmod{15}$ \checkmark
\begin{ergebnis} $z \equiv 110 \pmod{255}$, also $z = 110 + 255k$ ($k \in \Z$); kleinste positive Lösung $z = 110$. \end{ergebnis}

% ---------------------------------------------------------------------------
\section*{Aufgabe 5: Modulo-11-Code $[10, 8, 3]_{11}$}
\teil{a} Kontrollmatrix (Stellen $1, \dots, 10$):
\[
H = \begin{pmatrix}
1 & 1 & 1 & 1 & 1 & 1 & 1 & 1 & 1 & 1\\
1 & 2 & 3 & 4 & 5 & 6 & 7 & 8 & 9 & 10
\end{pmatrix}
\]
Ein Wort $c$ ist Codewort genau dann, wenn $H c^T = 0$, also
$\sum_{i=1}^{10} c_i \equiv 0$ und $\sum_{i=1}^{10} i \cdot c_i \equiv 0 \pmod{11}$.

\teil{b} Kontrollziffern für $a = 1\,1\,5\,0\,0\,0\,0\,5$ ($c_i = a_i$ für $i = 1, \dots, 8$):
\begin{align*}
c_9 &\equiv \sum_{i=1}^{8} (1+i)\, c_i = 2 \cdot 1 + 3 \cdot 1 + 4 \cdot 5 + 9 \cdot 5 = 70 \equiv 4 \pmod{11},\\
c_{10} &\equiv \sum_{i=1}^{8} (9-i)\, c_i = 8 \cdot 1 + 7 \cdot 1 + 6 \cdot 5 + 1 \cdot 5 = 50 \equiv 6 \pmod{11}.
\end{align*}
\begin{ergebnis} $c_9 = 4$, $c_{10} = 6$, $c = 1\,1\,5\,0\,0\,0\,0\,5\,4\,6$ \end{ergebnis}

\teil{c} Syndrome $S(r) = H r^T = (s_1, s_2)^T$:
\begin{align*}
r_1 = 1\,1\,5\,0\,0\,0\,0\,5\,1\,1:\quad
s_1 &= 1 + 1 + 5 + 5 + 1 + 1 = 14 \equiv 3,\\
s_2 &= 1 + 2 + 15 + 40 + 9 + 10 = 77 \equiv 0 \quad\Rightarrow\quad S(r_1) = (3, 0)^T,\\[4pt]
r_2 = 1\,1\,5\,0\,0\,0\,0\,7\,3\,3:\quad
s_1 &= 1 + 1 + 5 + 7 + 3 + 3 = 20 \equiv 9,\\
s_2 &= 1 + 2 + 15 + 56 + 27 + 30 = 131 \equiv 10 \quad\Rightarrow\quad S(r_2) = (9, 10)^T .
\end{align*}

\teil{d} Bei genau einem Fehler der Größe $e_x$ an der Stelle $x$ gilt $S = e_x \cdot (1, x)^T$,
also $e_x = s_1$ und $x = s_2 / s_1$.

Für $r_1$: $x = 0 / 3 = 0$. Die Stelle $0$ gibt es nicht (Stellen $1, \dots, 10$).
Das Syndrom passt also zu keinem Einzelfehler: $r_1$ enthält mindestens zwei Fehler.
Der Code hat $d = 3$ und korrigiert nur $1$ Fehler.
\begin{ergebnis} $r_1$ ist nicht dekodierbar (mindestens 2 Fehler). \end{ergebnis}

\teil{e} Für $r_2$: $e_x = s_1 = 9$ und $x = s_2 / s_1 = 10 \cdot 9^{-1} = 10 \cdot 5 = 50 \equiv 6 \pmod{11}$
(denn $9 \cdot 5 = 45 \equiv 1$). Fehler an Stelle $6$ mit Größe $9$:
\[
c_6 = r_6 - 9 = 0 - 9 \equiv 2 \pmod{11}.
\]
\begin{ergebnis} $c = 1\,1\,5\,0\,0\,2\,0\,7\,3\,3$ \end{ergebnis}

\teil{f} Probe $H c^T = 0$:
\begin{align*}
\text{Zeile 1:}\quad & 1 + 1 + 5 + 2 + 7 + 3 + 3 = 22 \equiv 0 \pmod{11},\\
\text{Zeile 2:}\quad & 1 + 2 + 15 + 12 + 56 + 27 + 30 = 143 = 13 \cdot 11 \equiv 0 \pmod{11}.
\end{align*}
Also ist $c$ ein Codewort. \checkmark

\medskip
{\small\textit{Hinweis:} Andere Formen von $H$ (z.\,B. systematisch $H = (A \mid E)$) beschreiben denselben Code.
Die Syndromwerte sehen dann anders aus, das Ergebnis (Fehlerstelle, Codewort) ist dasselbe.}

\end{document}
